\section{物理层 II：分岔、阈值与临界行为}

\begin{definition}[分岔式变化]
\textbf{分岔式变化}是指某个参数变化改变了动力学的定性结构，而不只是改变了运动速度的数值大小。
\end{definition}

分岔语言之所以重要，是因为许多干预的价值恰恰在于：
它们在努力成本上很小，但在结构后果上很大。
一个新的闹钟、一条改变后的路线、一个被屏蔽的应用，或一条预写好的任务提示，
在能量消耗上看似很小，但在结果上可能很大，
因为它已经把某个控制参数推过了阈值。

\begin{example}[通勤到健身房的路径设计]
如果一位上班族先回家，晚上就可能滑入休息盆地。
如果路线被改为通勤途中直接经过健身房入口，系统就可能根本不会重新构造旧盆地。
这个干预在描述上很小，但在后果上很大，
因为它改变了决策发生的时间与地点。
这正是分岔语言旨在捕捉的那类对机制敏感的推理。
\end{example}